\section{Conclusion}
\label{sec:conclusion}

We asked what an embodied language-model agent does with an instruction it
was given and then told to drop, in the situation where nobody is asking:
the cue appears in the world, out of view, while the agent is busy. Across
six model families, two simulators and six memory designs, the answer
depends almost entirely on how the intention is stored. The stores in
common use, a re-injected briefing, retrieval or a rolling summary,
produce the action after cancellation in 12 to 64\% of episodes. An entry
with a status, rewritten when the user withdraws the instruction, brings
this to 7\% or less on every model at no systematic cost to execution. A model
can keep that entry itself, so the design is a method rather than an upper
bound, though the two smallest bookkeepers invent obligations that were never given,
and with six standing intentions the advantage is intact for two of the
four actors we tested, reduced for one and gone for one. Retrospective memory asks whether an
agent can recall what it was told; an agent that acts on its own must also
be able to drop it.
